% Part: lambda-calculus
% Chapter: introduction
% Section: lambda-computable

\documentclass[../../../include/open-logic-section]{subfiles}

\begin{document}

\olfileid{lam}{int}{cmp}
\olsection{الدوال \usetoken{S}{lambda definable} قابلة للحساب}

\begin{thm}
\ollabel{thm:lambda-computable}
كل دالة جزئية~${\OLId{latin-ordinary}{f}}$ تكون !!{lambda defined} بحد لامبدا تقبل الحساب.
\end{thm}

\begin{proof}
لتكن الدالة~${\OLId{latin-ordinary}{f}}$ !!{lambda defined} بحد لامبدا~${\OLId{latin-looped}{X}}$.
نصف أولًا إجراء حساب~${\OLId{latin-ordinary}{f}}$ وصفًا غير صوري.
إذا أعطينا~${\OLId{latin-ordinary}{m}}_{\OLMathNumeral{0}}$، \dots،~${\OLId{latin-ordinary}{m}}_{{\OLId{latin-ordinary}{n}}-{\OLMathNumeral{1}}}$،
كتبنا الحد~${\OLId{latin-looped}{X}} \num {\OLId{latin-ordinary}{m}}_{\OLMathNumeral{0}} \ldots \num {\OLId{latin-ordinary}{m}}_{{\OLId{latin-ordinary}{n}}-{\OLMathNumeral{1}}}$.
ثم أنشأنا شجرة: في طبقتها الأولى جميع اختزالات هذا الحد
بخطوة واحدة، وتحت كل حد منها جميع اختزالاته بخطوة واحدة،
فتكون الطبقة الثانية اختزالات الحد الأصلي بخطوتين،
وهكذا طبقة بعد طبقة. متى بلغنا تمثيلًا عدديًا أعدناه جوابًا؛
وإن لم نبلغه كانت الدالة غير معرّفة.

وبأطروحة تشيرش نستنتج من هذا الوصف قابلية الدالة للحساب.
والأدق أن نصف البحث وصفًا عوديًا. فلنا مثلًا أن نعرّف دوال
وعلاقات عودية بدائية باسم
«$\fn{IsASubterm}$» و«$\fn{Substitute}$»
و«$\fn{ReducesToInOneStep}$» و«$\fn{ReductionSequence}$»
و«$\fn{Numeral}$»، ونحوها.
ثم يكون حساب~${\OLId{latin-ordinary}{f}}({\OLId{latin-ordinary}{m}}_{\OLMathNumeral{0}}, \dots, {\OLId{latin-ordinary}{m}}_{{\OLId{latin-ordinary}{n}}-{\OLMathNumeral{1}}})$ إجراءً عوديًا جزئيًا:
نبحث عن متتالية اختزالات، كل منها بخطوة واحدة، تبدأ
بـ${\OLId{latin-looped}{X}} \num{{\OLId{latin-ordinary}{m}}_{\OLMathNumeral{0}}} \dots \num{{\OLId{latin-ordinary}{m}}_{{\OLId{latin-ordinary}{n}}-{\OLMathNumeral{1}}}}$ وتنتهي بتمثيل عددي،
فنعيد العدد الطبيعي الذي يمثله.
وتفصيل ذلك طويل مضجر، إلا أنه جار على الطرق المعتادة.
\end{proof}

\end{document}
